\section{הרצאה 13: משפט שטורם-ליוביל}
\date{13.11.2025}
\label{sec:lecture13}
\begin{mainbox}
    {מידע כללי} \textbf{תאריך:} 13 בנובמבר 2025\\
    \textbf{נושא:} אופרטורים דיפרנציאליים, צמודים לעצמם, משפט שטורם-ליוביל והכללת
    טורי פורייה.
\end{mainbox}
\subsection*{שאלות מנחות}
\begin{enumerate}[label=\textbf{\arabic{*}.}, rightmargin=2em, leftmargin=1em]
    \item כיצד מגדירים אופרטור צמוד (Adjoint) לאופרטור דיפרנציאלי?

    \item מהו התנאי לכך שאופרטור יהיה צמוד לעצמו (Self-Adjoint)?

    \item מהי צורת שטורם-ליוביל ומדוע היא חשובה?

    \item מהם שלושת החלקים העיקריים של משפט שטורם-ליוביל?

    \item כיצד מוכיחים שהערכים העצמיים ממשיים ושפונקציות עצמיות אורתוגונליות?

    \item כיצד פורסים פונקציה כללית באמצעות הפונקציות העצמיות (טור פורייה מוכלל)?
\end{enumerate}
\subsection*{רעיונות מרכזיים והגדרות}
\subsubsection*{אופרטורים דיפרנציאליים לינאריים}
נתייחס לאופרטור דיפרנציאלי לינארי מסדר שני $L$ הפועל על פונקציה $u(x)$ בקטע
$[a,b]$:
\[
    L[u] = P_{0}(x) u'' + P_{1}(x) u' + P_{2}(x) u \]ניתן לחשוב על האופרטור $L$ כ
    אנלוגיה אינסוף-ממדית למטריצה הפועלת על וקטור.
    \begin{definitionbox}
        {האופרטור הצמוד (Adjoint Operator)} נגדיר את האופרטור הצמוד $\bar{L}$ (או
        $L^{\dagger}$) באופן הבא:

        \[
            \bar{L}[u] = \frac{d^{2}}{dx^{2}}(P_{0} u) - \frac{d}{dx}(P_{1} u) +
            P_{2} u
        \]שימו לב לשינוי בסימנים ולסדר הגזירה (הנגזרות פועלות על המכפלה של המקדם
        והפונקציה).
    \end{definitionbox}
     \subsubsection*{אופרטור צמוד לעצמו (Self-Adjoint)}
מחפשים את התנאים בהם האופרטור שווה לצמוד של עצמו, כלומר $L[u] = \bar{L}[u]$. נפתח את הביטוי של האופרטור הצמוד:
\begin{align*}
\bar{L}[u] &= (P_0 u)'' - (P_1 u)' + P_2 u \\
&= (P_0' u + P_0 u')' - (P_1' u + P_1 u') + P_2 u \\
&= P_0'' u + 2P_0' u' + P_0 u'' - P_1' u - P_1 u' + P_2 u \\
&= P_0 u'' + (2P_0' - P_1)u' + (P_0'' - P_1' + P_2)u
\end{align*}
נשווה מקדמים עם $L[u] = P_{0} u'' + P_{1} u' + P_{2} u$:
\begin{itemize}
\item מקדם $u''$: $P_{0} = P_{0}$ (מתקיים תמיד).
\item מקדם $u'$: $P_{1} = 2P_{0}' - P_{1} \implies 2P_{1} = 2P_{0}' \implies \boxed{P_1 = P_0'}$.
\end{itemize}
    אם התנאי $P_{1} = P_{0}'$ מתקיים, המקדם של $u$ במשוואת הצמוד הופך ל- $P_{0}''
    - P_{0}'' + P_{2} = P_{2}$, וזה תואם ל-$L$ המקורי.
    \begin{resultbox}
        {צורת שטורם-ליוביל} כאשר האופרטור צמוד לעצמו ($P_{1} = P_{0}'$), ניתן
        לכתוב אותו בצורה קומפקטית:

        \[
            L[u] = \frac{d}{dx}\left(P(x)\frac{du}{dx}\right) + Q(x)u
        \]כאשר סימנו $P(x) \equiv P_{0}(x)$ ו-$Q(x) \equiv P_{2}(x)$.
    \end{resultbox}
\subsubsection*{בעיית שטורם-ליוביל}
משוואת שטורם-ליוביל מוגדרת כך:
\[
    L[u] + \lambda w(x) u = 0
\]
או בצורה מפורשת:
\[
    \boxed{\frac{d}{dx}\left(P(x)\frac{du}{dx}\right) + Q(x)u + \lambda w(x)u = 0}
\]
כאשר:
\begin{itemize}
    \item $\lambda$ הוא פרמטר ספקטרלי (ערך עצמי).
    \item $w(x)$ היא \textbf{פונקציית משקל} (Weight Function), $w(x) > 0$.
\end{itemize}
זוהי בעיית ערכים עצמיים מוכללת: $L[u] = -\lambda w u$. באנלוגיה למטריצות: $Av = \lambda v$.  

\subsection*{דוגמה מעובדת: מיתר קשור בקצותיו}
\begin{examplebox}
    {פתרון בעיית שטורם-ליוביל פשוטה} נפתור את המשוואה בקטע $[0, L]$:
    \[
        u'' + \lambda u = 0
    \]
    עם תנאי שפה הומוגניים: $u(0) = 0, \quad u(L) = 0$.
    כאן: $P(x)=1, Q(x)=0, w(x)=1$.
    \textbf{מקרה 1: $\lambda > 0$} נסמן $\lambda = k^{2}$. הפתרון הכללי:
    \[
        u(x) = A \cos(kx) + B \sin(kx)
    \]
    תנאי שפה:
    \begin{itemize}
        \item $u(0) = A \cdot 1 + 0 = 0 \implies A=0$.
        \item $u(L) = B \sin(kL) = 0$.
    \end{itemize}
    עבור פתרון לא טריוויאלי ($B \neq 0$), נדרש $\sin(kL) = 0$, כלומר $kL = n\pi$ עבור $n=1,2,\dots$.
    קיבלנו ערכים עצמיים ופונקציות עצמיות:
    \[
        \lambda_{n} = \left(\frac{n\pi}{L}\right)^{2}, \quad u_{n}(x) = \sin\left(\frac{n\pi x}{L}\right)
    \]
    \textbf{מקרה 2: $\lambda < 0$} נסמן $\lambda = -\mu^{2}$. הפתרון:
    \[
        u(x) = A e^{\mu x} + B e^{-\mu x}\quad \text{או}\quad C \sinh(\mu x) + D \cosh(\mu x)
    \]
    בדיקת תנאי שפה מראה שרק הפתרון הטריוויאלי $u(x) \equiv 0$ מתקיים (סינוס היפרבולי מתאפס רק ב-0).
    \textbf{מקרה 3: $\lambda = 0$} המשוואה $u''=0 \implies u(x) = Ax+B$.
    תנאי השפה מחייבים $A=B=0$.
    \textbf{מסקנה:} קיימים אינסוף ערכים עצמיים בדידים $\lambda_{n} > 0$, והפונקציות העצמיות הן סינוסים.
\end{examplebox}

\subsection*{משפט שטורם-ליוביל}
בהינתן אופרטור $L$ צמוד לעצמו ותנאי שפה מתאימים (כגון $u$ או $u'$ מתאפסים בקצוות, או תנאים מחזוריים), מתקיימות התכונות הבאות:  
\begin{theorembox}{משפט שטורם-ליוביל}
    \begin{enumerate}
        \item \textbf{ערכים עצמיים ממשיים:} כל הערכים העצמיים $\lambda$ הם מספרים ממשיים (גם אם הפונקציות העצמיות מרוכבות).
        \item \textbf{אורתוגונליות:} פונקציות עצמיות המתאימות לערכים עצמיים שונים הן אורתוגונליות ביחס לפונקציית המשקל $w(x)$:
        \[
            \text{אם }\lambda_{n} \neq \lambda_{m} \implies \langle u_{n}, u_{m} \rangle_{w} \equiv \int_{a}^{b} u_{n}^{*}(x) u_{m}(x) w(x)  dx = 0
        \]
        \item \textbf{שלמות (Completeness):} אוסף הפונקציות העצמיות $\{u_{n}\}_{n=1}^{\infty}$ מהווה בסיס שלם במרחב הפונקציות. כלומר, כל פונקציה "סבירה" $f(x)$ ניתנת לפריסה כטור:
        \[
            f(x) = \sum_{n=1}^{\infty} a_{n} u_{n}(x)
        \]
    \end{enumerate}
\end{theorembox}

\subsection*{הוכחות}
\subsubsection*{זהות לגראנז' והרמיטיות}
נגדיר מכפלה פנימית:
\[
    \langle u, v \rangle = \int_{a}^{b} u^{*}(x)\, v(x)\, dx
\]
(כעת בלי משקל). נבדוק מתי האופרטור הוא \textbf{הרמיטי} (Hermitian), כלומר מתי
$\langle u, L v \rangle = \langle L u, v \rangle$. עבור אופרטור צמוד לעצמו מהצורה
$L[u]=(P u')' + Q u$ נחשב את ההפרש האינטגרלי:
\begin{align*}
\int_{a}^{b} \big(u^{*} L v - (L u)^{*} v\big)\,dx
&= \int_{a}^{b} \Big[ u^{*}\big((P v')' + Q v\big) - \big((P u')' + Q u\big)^{*} v \Big]\,dx \\
&= \int_{a}^{b} \Big[ u^{*}(P v')' - \big((P u')'\big)^{*} v \Big]\,dx
\qquad (\text{נניח }P,Q\text{ ממשיים so }Q^{*}=Q)\\
&= \int_{a}^{b} \frac{d}{dx}\!\big[ P(x)\big(u^{*} v' - u'^{*} v\big)\big]\,dx \\
&= \big[\,P(x)\big(u^{*} v' - u'^{*} v\big)\,\big]_{x=a}^{x=b}.
\end{align*}
הביטוי בגבולות הוא איבר השפה הכללי (הוורונסקיאן המוכלל). אם תנאי השפה מבטלים את
הביטוי הזה (למשל $u(a)=u(b)=0$, או תנאים המחייבים את הביטוי בגבולות להתאפס,
או תנאים מחזוריים), אזי ההפרש מתאפס ולפיכך האופרטור הרמיטי:
\[
    \langle u, L v \rangle = \langle L u, v \rangle .
\]

\subsubsection*{הוכחת ערכים עצמיים ממשיים}
נניח $L u_{n} = -\lambda_{n} w u_{n}$. נשתמש בהרמיטיות עם $v=u_{n}$:
\[
    \langle u_{n}, L u_{n} \rangle = \langle L u_{n}, u_{n} \rangle
\]
נציב את המשוואה:
\[
    \langle u_{n}, -\lambda_{n} w u_{n} \rangle = \langle -\lambda_{n} w u_{n}, u_{n} \rangle
\]
נוציא את הקבועים מהאינטגרל:
\[
    -\lambda_{n} \int_a^b u_{n}^{*} u_{n} w  dx = -\lambda_{n}^{*} \int_a^b u_{n}^{*} u_{n} w  dx
\]
מכיוון שהאינטגרל $\int_a^b |u_{n}|^{2} w  dx$ חיובי ואינו אפס, נקבל:
\[
    \lambda_{n} = \lambda_{n}^{*}
\]
לכן $\lambda_{n}$ ממשי.

\subsubsection*{הוכחת אורתוגונליות}
נניח שני פתרונות עם ערכים עצמיים שונים:
\begin{enumerate}
    \item $L u_{n} + \lambda_{n} w u_{n} = 0$
    \item $L u_{m} + \lambda_{m} w u_{m} = 0$
\end{enumerate}
נכפיל את (1) ב-$u_{m}^{*}$ ואת הצמוד של (2) ב-$u_{n}$, ונחסר ביניהם. לאחר אינטגרציה ושימוש בהרמיטיות האופרטור (אגף שמאל מתאפס), נקבל:
\[
    (\lambda_{n} - \lambda_{m}) \int_{a}^{b} u_{m}^{*} u_{n} w  dx = 0
\]
מכיוון ש-$\lambda_{n} \neq \lambda_{m}$, האינטגרל חייב להתאפס:
\[
    \langle u_{m}, u_{n} \rangle_{w} = 0
\]

\subsection*{טור פורייה מוכלל}
מכיוון שהפונקציות העצמיות מהוות בסיס שלם, ניתן לפרוס כל פונקציה $f(x)$:
\[
    f(x) = \sum_{n=1}^{\infty} a_{n} u_{n}(x)
\]
כדי למצוא את המקדמים $a_{n}$, נכפיל את שני האגפים ב-$u_{m}^{*}(x) w(x)$ ונבצע אינטגרציה בקטע $[a,b]$:
\[
    \int_{a}^{b} f(x) u_{m}^{*}(x) w(x)  dx = \sum_{n=1}^{\infty} a_{n} \int_{a}^{b} u_{n}(x) u_{m}^{*}(x) w(x)  dx
\]
בגלל האורתוגונליות, כל האיברים בסכום מתאפסים למעט $n=m$. נקבל:
\[
    \langle u_{m}, f \rangle_{w} = a_{m} \langle u_{m}, u_{m} \rangle_{w}
\]
ולכן הנוסחה למקדם היא:
\[
    \boxed{a_n = \frac{\int_{a}^{b} f(x) u_{n}^{*}(x) w(x)  dx}{\int_{a}^{b} |u_{n}(x)|^{2} w(x)  dx}}
\]
אם הפונקציות מנורמלות ($\int_a^b |u_{n}|^{2} w  dx = 1$), המכנה הוא 1.

\begin{summarybox}
    {סיכום}
    \begin{itemize}
        \item אופרטור $L$ הוא צמוד לעצמו אם $P_{1} = P_{0}'$.
        \item משוואת שטורם-ליוביל: $(Pu')' + Qu + \lambda w u = 0$.
        \item תחת תנאי שפה מתאימים, האופרטור הרמיטי.
        \item הערכים העצמיים ממשיים, והפונקציות העצמיות אורתוגונליות ביחס למשקל $w$.
        \item ניתן לפרוס פונקציות כטור של הפונקציות העצמיות (הכללה של טורי פורייה).
        \item בדוגמה הקלאסית ($u''+\lambda u=0$), קיבלנו את הבסיס של סינוסים (טורי פורייה סינוס).
    \end{itemize}
\end{summarybox}